\section{Background}

This section provides necessary background on zero-trust architecture, AI agents in HPC, HPC security, and prompt injection attacks.

\subsection{Zero-Trust Architecture}

Zero-Trust Architecture (ZTA), formalized in NIST SP 800-207~\cite{rose2020zero}, operates on the principle of ``never trust, always verify'', treating every access request as potentially hostile regardless of origin. Core tenets include per-session least-privilege access, dynamic decisions informed by multiple signals, and continuous monitoring. While widely adopted in enterprise and cloud environments, applying ZTA to HPC presents unique challenges: performance sensitivity of scientific workloads, shared-resource clusters including filesystems, interconnects, and schedulers, and the need to preserve collaborative, low-friction access patterns.

Recent work has begun addressing this gap. Alam et al.~\cite{alam2024federated} deployed federated SSO with zero-trust controls for the Isambard-AI/HPC infrastructures in the UK. Duckworth et al.~\cite{duckworth2023spiffe} proposed SPIFFE/SPIRE for workload identity in HPC~\cite{spiffe2024secure}. Macauley and Bhasker~\cite{macaulay2024high} measured ZTA maturity implementation effort in HPC, finding the identity pillar particularly challenging due to cost and complexity. Our work extends ZTA to a dimension these efforts do not address: \textit{behavioral trust of autonomous AI agents}. Existing ZTA in HPC focuses on identity verification; for AI agents, identity verification must be complemented by behavioral verification, attesting not just to who the agent is, but to what it does.

\subsection{AI Agents in HPC}

AI agents, autonomous systems that perceive, reason, and act to achieve goals, are entering HPC workflows~\cite{wooldridge2009introduction}, making dynamic decisions about data analysis, simulation parameters, and tool invocation during execution. Frameworks like Academy~\cite{kamatar2025empowering} and RHAPSODY~\cite{alsaadi2025rhapsody} now support agent deployment across federated HPC ecosystems, while AgentBound~\cite{buhler2025securing} addresses access control for MCP servers, the emerging standard for connecting agents to external tools~\cite{hou2025model}. Agent safety evaluation benchmarks~\cite{zhang2024agent} reveal significant vulnerabilities in current tool-use patterns.

HPC environments introduce unique characteristics for agent operation: shared parallel filesystems where access is governed by user-level POSIX permissions rather than project boundaries; multi-tenant compute nodes where schedulers like Slurm~\cite{yoo2003slurm} co-locate jobs from different users sharing kernel-level resources, including \textit{/tmp} and shared memory; data-intensive workflows where agents routinely process terabytes of untrusted scientific data as authoritative input; and API-driven intelligence where LLM backends accessed over HTTPS create encrypted, whitelisted exfiltration channels by design.

\subsection{Security in HPC}

HPC security relies on a perimeter model: Kerberos/SSH authentication~\cite{neuman1994kerberos} and federated identity through CILogon~\cite{basney2013cilogon}, RBAC through schedulers and POSIX permissions, boundary firewalls with internal traffic largely unencrypted for performance, and job accounting logs via slurmdbd~\cite{yoo2003slurm}. User-based firewalls enhance HPC security without disrupting workflows. This model has well-documented gaps: lateral movement once authenticated, credential theft granting full access, overprovisioned filesystem permissions, and critically, no behavioral monitoring where the system verifies identity but not intent. These gaps are tolerable for human users constrained by awareness; they become critical for AI agents that follow instructions blindly, including adversarial ones. An agent has valid credentials (passes authentication), legitimate permissions (passes authorization), and consistent behavior patterns (passes analytics), yet executes an attacker's commands.

\subsection{Prompt Injection and Agent Security}

Prompt injection, subverting an agent's instruction-following through adversarial inputs~\cite{he2025security}, has emerged as a fundamental security challenge for LLM-based systems. Prior work focuses on web-based scenarios where agents encounter malicious content on the internet~\cite{greshake2023not}. Existing defenses including input sanitization, instruction hierarchy, and output filtering have fundamental limitations: sanitization is incomplete due to unbounded payload space, instruction hierarchy is fragile as demonstrated by universal adversarial suffix attacks bypassing LLM guardrails~\cite{zou2023universal}, and output filtering is probabilistic since ML classifiers have inherent error rates. Comprehensive surveys~\cite{barcha2026systematic} catalogued defenses across 88+ studies, proposing expanded taxonomies for adversarial ML defenses. Critically, this literature has not addressed HPC contexts where injection exploits shared infrastructure rather than web content.